\documentclass[journal,twoside]{IEEEtran}
\usepackage[T1]{fontenc}
\usepackage{graphicx}
\usepackage{booktabs}
\usepackage{multicol,multirow}
\usepackage{newtxmath}
\usepackage{mathtools}
\usepackage{bm}
\usepackage[caption=false,font=footnotesize,subrefformat=parens]{subfig}
\usepackage{lipsum}
\usepackage{tabularx}
\usepackage[dvipsnames,table]{xcolor}
\usepackage{my-tikz}
\usepackage{colortbl}
\usepackage{scalerel}
\usepackage{threeparttable}
\usepackage{algorithm}
\usepackage[italicComments=false,commentColor=darkgray,noEnd=true]{algpseudocodex}
\usepackage{upgreek}
\usepackage{calc}
\usepackage[noadjust]{cite}
\usepackage{circuitikz}
\usepackage{url}
\usepackage[colorlinks,urlcolor=black,linkcolor=blue,citecolor=blue]{hyperref}

\usetikzlibrary{svg.path}
\definecolor{orcidlogocol}{HTML}{A6CE39}
\tikzset{
  orcidlogo/.pic={
    \fill[orcidlogocol] svg{M256,128c0,70.7-57.3,128-128,128C57.3,256,0,198.7,0,128C0,57.3,57.3,0,128,0C198.7,0,256,57.3,256,128z};
    \fill[white] svg{M86.3,186.2H70.9V79.1h15.4v48.4V186.2z}
                 svg{M108.9,79.1h41.6c39.6,0,57,28.3,57,53.6c0,27.5-21.5,53.6-56.8,53.6h-41.8V79.1z M124.3,172.4h24.5c34.9,0,42.9-26.5,42.9-39.7c0-21.5-13.7-39.7-43.7-39.7h-23.7V172.4z}
                 svg{M88.7,56.8c0,5.5-4.5,10.1-10.1,10.1c-5.6,0-10.1-4.6-10.1-10.1c0-5.6,4.5-10.1,10.1-10.1C84.2,46.7,88.7,51.3,88.7,56.8z};
  }
}

\newcommand\orcidicon[1]{\href{https://orcid.org/#1}{\mbox{\scalerel*{
  \begin{tikzpicture}[yscale=-1,transform shape]
    \pic{orcidlogo};
  \end{tikzpicture}
}{|}}}}

% \makesavenoteenv{tabular}
\interdisplaylinepenalty=2500
\aboverulesep = 0mm \belowrulesep = 0mm
\newcolumntype{Y}{>{\centering\arraybackslash}X}
\newcolumntype{M}[1]{>{\centering\arraybackslash}m{#1}} % centered m
\newcommand{\tabvertspace}{\specialrule{0em}{0.08em}{.08em}}

\definecolor{doicolor}{RGB}{2, 48, 129}

\DeclareMathOperator*{\argmax}{\arg\max}
\DeclareMathOperator*{\argmin}{\arg\min}
\DeclareMathOperator*{\argmint}{\widetilde{\arg\min}}
\DeclareMathOperator*{\uargmin}{\mathrm{u}\arg\min}

\markboth{Southeast University Software-Defined Radio Experiment Course, January 2024}{W. Zhao: Dual-Mode PSK Transceiver on SDR With FPGA}
\IEEEpubid{
  \begin{tabular}[t]{c}
    \copyright{} 2024 Wuqiong Zhao.
    This work is licensed under a Creative Commons Attribution-ShareAlike 4.0 License.\\
    For more information, see \url{https://creativecommons.org/licenses/by-sa/4.0/}.
  \end{tabular}
}

\hypersetup{
  , pdftitle    = {Dual-Mode {PSK} Transceiver on {SDR} With {FPGA}}
  , pdfauthor   = {Wuqiong Zhao}
  , pdfsubject  = {Southeast University Software-Defined Radio Experiment Course}
  , pdfkeywords = {Phase-shift keying (PSK), software-defined radio (SDR), transceiver design, modulation, demodulation, field programmable gate array (FPGA).}
}

\begin{document}

  \title{Dual-Mode PSK Transceiver on SDR With FPGA}

  \author{%
    Wuqiong~Zhao{\hspace{.1em}\textsuperscript{\orcidicon{0000-0002-9550-7423}}},~\IEEEmembership{Student Member, IEEE}
    \thanks{Date of publication 1 January 2024; date of current version 1 January 2024.
      The author would like to thank Lecturer Dan Yang for the guidance and help in the course of the experiment.}
    \thanks{Wuqiong Zhao is with Southeast University, Nanjing 211189, China. (e-mail: wqzhao@seu.edu.cn).}
    \thanks{Online URL: \url{https://go.wqzhao.org/sdr-psk-fpga}}
    % \vspace{-1mm}
  }

  \maketitle

  \begin{abstract}
    In this experiment, we implement a dual-mode PSK transceiver on SDR with FPGA,
    supporting both BPSK and QPSK.
    Moreover, the transceiver is designed to be able to switch between the two modes by introducing packet-based communication,
    where modulation information can be extracted from the packet header.
  \end{abstract}
  \begin{IEEEkeywords}
    Phase-shift keying (PSK), software-defined radio (SDR), transceiver design, modulation, demodulation, field programmable gate array (FPGA).
  \end{IEEEkeywords}

  \section{Introduction}

    \IEEEPARstart{S}{oftware-defined} radio is interesting.
    Channel estimation is also interesting!
    \cite{zhao2023ompl,you2023beam}.

  \section{System Overview}

    \subsection{Software-Defined Radio}

    \subsection{Transceiver Design}

    \subsection{BPSK/QPSK Modulation}

  \section{Transmitter}

  \section{Receiver}

    \subsection{Overview}

    \subsection{Carrier Synchronization Using Costas Loop}

    \subsection{Symbol Synchronization Using Gardner Loop}

  \section{Packet-Based Communication}

    \subsection{Frame Structure}

    \subsection{Packetizer Design}

    \subsection{Depacketizer Design}

    \subsection{SPB Detection}

      \subsubsection{Strength Detection (SD)}

      \subsubsection{Packet Detection (SD)}

      \subsubsection{Boundary Detection (BD)}

  \section{Experiment Results}

  \section{Discussions}

  \section{Conclusion}

  \bibliographystyle{IEEEtran}
  \bibliography{IEEEabrv, ref}

\end{document}
